% !TEX root = ../jlt.tex
\subsection{Translation into System F}
\label{app:fjlt:translation}

The target is System~F with products, sums, the carriers of $\Sigma$ taken as constants with their total operations, and inductive types with their iterators: the natural numbers with $\mathit{iter} : \forall\sigma.\,\mathbb{N} \to (\sigma \to \sigma) \to \sigma \to \sigma$, reducing by $\mathit{iter}\,(k{+}1)\,s\,z \longrightarrow \mathit{iter}\,k\,s\,(s\,z)$ and $\mathit{iter}\,0\,s\,z \longrightarrow z$, and for each term type $\mu F$ its iterator $\mathit{fold}_F : \forall\sigma.\,(F(\sigma) \to \sigma) \to \mu F \to \sigma$, reducing one constructor at a time, $\mathit{fold}_F\,\beta\,(\mathrm{in}\,n) \longrightarrow \beta\,(F(\mathit{fold}_F\,\beta)\,n)$; the paramorphism, which also hands each node itself to the algebra, is the iterator at $\mu F \times \sigma$ with the node rebuilt beside the result, as usual.
Reduction in the target is full $\beta$-reduction, any redex anywhere, together with projections, case, the $\delta$-rules of the constants, the unfolding of a defined name to its body and the two iterator rules; no evaluation order is imposed, which is the setting of the normalisation and confluence results used below.
Records are nested pairs.
System~F with products, sums and positive inductive types with their iterators is strongly normalising and confluent, and its iterative fragment embeds into System~F by an encoding that preserves reduction~\cite{matthes1999extensions}; for first-order rewrite rules on algebraic constants, Breazu-Tannen and Gallier show that adding them to System~F conserves strong normalisation and confluence~\cite{breazutannen1991polymorphic,breazutannen1994polymorphic}.
These results cover the parts separately, and this appendix assumes their combination, which the same method gives: $\mathit{iter}$ in the accumulator form above is the standard iterator at $\sigma \to \sigma$, each of its steps at least one step of that; the defined names are acyclic, by the declaration order, and expand away; the $\delta$-rules are ground, a constant out of constants, and preserve reducibility; and confluence then follows by Newman's lemma, there being no critical pairs.
That is a sketch, not a proof, and the one-constructor-at-a-time iterators are what make Theorem~\ref{thm:simulation} exact where a Church encoding, unfolding all at once, would not be.
Figure~\ref{fig:fjlt-translation} gives the translation.

\paragraph{Reading the translation.}
A function object is a pair of its fields and its call function, so a fieldless \cppinline{std::plus<size_t>} is $((), \mathit{plus})$, passing it as a step uses its second component, and reading a field uses its first.
A trait has no image of its own; its specialisations are closed terms $f_{\overline{\tau}}$, and a generic function that serves as an instance with a context is the function from records to a record that \textsc{E-Shape} builds its dictionary with.
Where \cppinline{size_t} is read as $\mathbb{N}$, by a coercion constant from the carrier to the inductive naturals, $\llbracket n \rrbracket$ is the natural number that is $n$'s value; the translation of a term with a runtime \cppinline{n} is the term with $n$ as a free variable of type $\mathbb{N}$, and Theorem~\ref{thm:simulation} is stated for closed terms.
The call function of a constrained operation abstracts over its dictionaries as a generic function does, and an application supplies them.
The iterator is the catamorphism of the chain of a natural number, a fold over any range of $\mathcal{R}$ is that catamorphism at the range's length with the range's position carried beside the state, and \cppinline{cata} is the catamorphism of a term's inductive type; a term type translates to $\mu F$, whose constructors are its node operations, and \textsc{T-Node} reads a construction as a term of that type, $\mathrm{in}_O$ of its fields, exactly where a term of type $\mu F$ is expected, the argument of a \cppinline{cata} or of an algebra, so that a term keeps one derivation.
The library's set recursion is the first two, and its one engine over terms is the third.
A box translates to a unary product rather than to nothing, on purpose: a box's only observable is its content, so which box holds it is erased, and the unit component keeps unboxing a step on both sides, which Lemma~\ref{lem:forward} needs.
Its one non-local ingredient is the \emph{dictionary} of an entailment: $\llbracket \Delta \Vdash Q \rrbracket$ is the record that witnesses $Q$, built by the same recursion as resolution.

\begin{figure}[htbp]
\scriptsize
\[
\resizebox{\linewidth}{!}{$\begin{array}{l@{\quad}l}
\multicolumn{2}{l}{\textit{types}} \\
\llbracket \alpha \rrbracket = \alpha \qquad \llbracket b \rrbracket = b & \\
\llbracket O\texttt{<}\overline{\tau}\texttt{>} \rrbracket = \overline{\llbracket \tau_y \rrbracket} \times (\forall\overline{\beta}.\,\overline{\llbracket Q \rrbracket} \to \overline{\llbracket \tau_x \rrbracket} \to \llbracket \tau \rrbracket) & \text{a function object: its fields and its call function} \\
\llbracket \texttt{std::pair<}\tau_1\texttt{, }\tau_2\texttt{>} \rrbracket = \llbracket \tau_1 \rrbracket \times \llbracket \tau_2 \rrbracket & \\
\llbracket \texttt{std::optional<}\tau\texttt{>} \rrbracket = 1 + \llbracket \tau \rrbracket \qquad \llbracket \texttt{std::variant<}\tau_1\texttt{, }\tau_2\texttt{>} \rrbracket = \llbracket \tau_1 \rrbracket + \llbracket \tau_2 \rrbracket & \\
\llbracket \texttt{std::shared\_ptr<const }\tau\texttt{>} \rrbracket = 1 \times \llbracket \tau \rrbracket & \text{a box is a unary product} \\
\llbracket R\texttt{<}\tau\texttt{>} \rrbracket = \llbracket R \rrbracket \qquad \llbracket \mu F \rrbracket = \mu F \text{ with } F\,X = \textstyle\sum_{O} \prod \overline{\llbracket \tau_y \rrbracket}[X / \text{the leg types}] & \text{a range's state; a term type, the sum of its nodes' fields, a leg's type replaced by } X \\[6pt]
\multicolumn{2}{l}{\textit{constraints, as record types}} \\
\llbracket C\texttt{<}\overline{\tau}\texttt{>} \rrbracket = \{\,\overline{r_i : \llbracket R_i \rrbracket[\overline{\tau}/\overline{\alpha}]}\,\} \text{ where } \texttt{concept } C \texttt{ = } \overline{R} & \text{one field per requirement} \\
\llbracket Q \rrbracket \text{ for a nested shape } Q & \text{its record type} \\
\llbracket \texttt{!}C\texttt{<}\overline{\tau}\texttt{>} \rrbracket = 1 & \text{a negation carries no evidence} \\
\llbracket e \rrbracket = \mathit{Bool} \text{ for a law } e & \\
\llbracket \texttt{\{ }e\texttt{ \} -> }S\texttt{<}\tau\texttt{>} \rrbracket = \overline{\llbracket \tau_x \rrbracket} \to \llbracket \tau \rrbracket & \text{an operation on parameters } \overline{x : \tau_x} \\[6pt]
\multicolumn{2}{l}{\textit{dictionaries, from entailment derivations}} \\
\llbracket \Delta \Vdash Q \rrbracket = q \text{ if } Q \in \Delta \text{ with record variable } q & \text{\textsc{E-Assume}} \\
\llbracket \Delta \Vdash C\texttt{<}\overline{\tau}\texttt{>} \rrbracket = \{\,\overline{r_i = \llbracket \Delta \Vdash R_i[\overline{\tau}/\overline{\alpha}] \rrbracket}\,\} & \text{\textsc{E-Shape}} \\
\llbracket \Delta \Vdash e \rrbracket = \mathit{true} & \text{\textsc{E-Law}, \textsc{E-Law-Assumed}} \\
\llbracket \Delta \Vdash \texttt{!}C\texttt{<}\overline{\tau}\texttt{>} \rrbracket = () & \text{\textsc{E-Not}} \\
\llbracket \Delta \Vdash \texttt{\{ }e_r\texttt{ \} -> }S\texttt{<}\tau\texttt{>} \rrbracket = \lambda\overline{a}.\,\llbracket \tau^{S}\texttt{(}e_r\texttt{)} \rrbracket & \text{\textsc{E-Op}, over the requirement's parameters, the conversion in the field} \\[6pt]
\multicolumn{2}{l}{\textit{declarations}} \\
\texttt{template <> constexpr } \tau\; f\texttt{<}\overline{\tau}\texttt{>(}\overline{x}\texttt{) ... } e \;\mapsto\; f_{\overline{\tau}} = \lambda\overline{x}.\,\llbracket e \rrbracket & \text{an instance} \\
\texttt{template <}\overline{\alpha}\texttt{> requires } \overline{Q}\texttt{ ... } f\texttt{(}\overline{x}\texttt{) ... } e \;\mapsto\; f = \Lambda\overline{\alpha}.\,\lambda\overline{q : \llbracket Q \rrbracket}.\,\lambda\overline{x}.\,\llbracket e \rrbracket & \text{a generic function} \\
\text{the same with a descending self-call on } p \in \overline{x} \;\mapsto\; f = \Lambda\overline{\alpha}.\,\lambda\overline{q}.\,\lambda\overline{x}.\,(\mathit{iter}\,(\mathit{toNat}\,p + 1)\,F\,d)\,\overline{x} & \text{the one exception, } F = \lambda f.\,\lambda\overline{x}.\,\llbracket e \rrbracket \\
\texttt{struct } O \texttt{ \{ } \overline{y}\texttt{; ... operator()(}\overline{x}\texttt{) ... } e \texttt{ ... \};} & \text{an operation's call function:} \\
\qquad \mapsto\; O_{\mathrm{call}} = \lambda\overline{y}.\,\Lambda\overline{\beta}.\,\lambda\overline{q : \llbracket Q \rrbracket}.\,\lambda\overline{x}.\,\llbracket e \rrbracket & \\
O \text{ standard} \;\mapsto\; O_{\mathrm{call}} = \lambda().\,o & \text{the standard's function } o \text{ as a constant} \\
\texttt{concept}, \texttt{= delete} \;\mapsto\; \text{nothing} & \\[6pt]
\multicolumn{2}{l}{\textit{terms}} \\
\llbracket f\texttt{<}\overline{\sigma}\texttt{>(}\overline{e}\texttt{)} \rrbracket = f_i\,\overline{\llbracket \sigma \rrbracket}\,\overline{\llbracket \Delta \Vdash Q[\overline{\sigma}/\overline{\alpha}] \rrbracket}\,\overline{\llbracket e \rrbracket} & \text{\textsc{T-Call}, $f_i$ the declaration it resolves to} \\
\llbracket f\texttt{<}\overline{\sigma}\texttt{>(}\overline{e}\texttt{)} \rrbracket = f_{\overline{\sigma}}\,\overline{\llbracket e \rrbracket} & \text{\textsc{T-Trait}, resolved to a specialisation} \\
\llbracket \tau^{S}\texttt{(}e_r[\overline{e}/\overline{a}]\texttt{)} \rrbracket = q_C.\pi.r\,\overline{\llbracket e \rrbracket} & \text{\textsc{T-Req-Assumed}, } \pi \text{ the path through nested shapes} \\
\llbracket O\texttt{<}\overline{\tau}\texttt{>\{}\overline{e_y}\texttt{\}} \rrbracket = (\overline{\llbracket e_y \rrbracket},\; O_{\mathrm{call}}\,\overline{\llbracket e_y \rrbracket}) & \text{construction; a fieldless } O \text{ has the empty tuple of fields} \\
\llbracket e\texttt{(}\overline{e'}\texttt{)} \rrbracket = (\pi_2\,\llbracket e \rrbracket)\,[\overline{\llbracket \sigma \rrbracket}]\,\overline{d}\,\overline{\llbracket e' \rrbracket} & \text{\textsc{T-Op}, \textsc{T-App}: application,} \\
\qquad \text{with } \overline{d} = \overline{\llbracket \Delta \Vdash Q[\overline{\tau}/\overline{\alpha}][\overline{\sigma}/\overline{\beta}] \rrbracket} \text{ the call operator's dictionaries} & \\
\llbracket e\texttt{.first} \rrbracket = \pi_1\,\llbracket e \rrbracket \qquad \llbracket e\texttt{.second} \rrbracket = \pi_2\,\llbracket e \rrbracket & \text{a pair's projections} \\
\llbracket e\texttt{.}y_i \rrbracket = \pi_i\,(\pi_1\,\llbracket e \rrbracket) & \text{a field, through the fields' pair} \\
\llbracket \tau\texttt{(}e\texttt{)} \rrbracket = \mathit{conv}_{\tau',\tau}\,\llbracket e \rrbracket & \text{a conversion, a constant of type } \llbracket \tau' \rrbracket \to \llbracket \tau \rrbracket \text{, the identity's } \delta\text{-step included at } \tau' = \tau \\
\llbracket f\texttt{<}\overline{\sigma}\texttt{>} \rrbracket^{\!\circ} = f\,\overline{\llbracket \sigma \rrbracket}\,\overline{d} & \text{the callable of a step: a function name,} \\
\llbracket e \rrbracket^{\!\circ} = (\pi_2\,\llbracket e \rrbracket)\,[\overline{\llbracket \sigma \rrbracket}]\,\overline{d} & \text{a function object; elaborated at the fold's types,} \\
\qquad \text{with } \overline{d} \text{ the dictionaries of its constraints at } \overline{\sigma} & \text{its dictionaries supplied} \\
\llbracket e \rrbracket^{\!\circ} = \lambda\overline{a}.\,q_C.\pi.r\,\overline{a} & \text{a step typed by \textsc{T-Req-Assumed}: the requirement's field} \\
\llbracket r \rrbracket = c_R\,\overline{\llbracket e \rrbracket} & \text{a range of instance } R \text{ with arguments } \overline{e} \\
\llbracket \texttt{size(}r\texttt{)} \rrbracket = |\llbracket r \rrbracket| \qquad \llbracket \texttt{head(}r\texttt{)} \rrbracket = \mathit{head}\,\llbracket r \rrbracket \qquad \llbracket \texttt{next(}r\texttt{)} \rrbracket = \mathit{next}\,\llbracket r \rrbracket & \text{a range's operations} \\
\llbracket \texttt{fold\_left(}r\texttt{, }a\texttt{, }f\texttt{)} \rrbracket = \pi_2\,(\mathit{iter}\,[\llbracket R \rrbracket \times \sigma]\,(\mathit{toNat}\,|\llbracket r \rrbracket|)\,\mathit{step}\,(\llbracket r \rrbracket, \llbracket a \rrbracket)) & \text{the iterator, over the range's length read as a natural} \\
\qquad \text{where } \mathit{step} = \lambda p.\,(\mathit{next}\,(\pi_1 p),\; \llbracket f \rrbracket^{\!\circ}\,(\pi_2 p)\,(\mathit{head}\,(\pi_1 p))) & \text{the position carried beside the state} \\
\qquad \text{with } |\cdot|, \mathit{head}, \mathit{next} \text{ the operations of } R & \\
\qquad \text{at a type variable } R \text{ (\textsc{T-Fold-R}): } \mathit{toNat} \circ q.\mathit{size}, q.\mathit{head}, q.\mathit{next} \text{ in their place} & q \text{ the record of } \mathit{IsRange}\texttt{<}R\texttt{, }\tau\texttt{>} \text{ in } \Delta \\
\llbracket \texttt{cata(}\mathit{alg}\texttt{, }t\texttt{)} \rrbracket = \pi_2\,(\mathit{fold}_F\,[\mu F \times \sigma]\,\beta\,\llbracket t \rrbracket) & \text{a fold over a term} \\
\qquad \text{where } \beta = \lambda n'.\,(\mathrm{in}\,(F\,\pi_1\,n'),\; \llbracket \mathit{alg} \rrbracket^{\!\circ}\,(\mathrm{in}\,(F\,\pi_1\,n'))\,(F\,\pi_2\,n')) & \\
\qquad & \text{the node rebuilt beside its result} \\
\llbracket t \rrbracket_{\mu F} = \mathit{node}_{O\texttt{<}\overline{\tau}\texttt{>}}\,\llbracket t \rrbracket & \text{\textsc{T-Node}: any term of node type read where } \mu F \text{ is expected} \\
\qquad \text{where } \mathit{node}_{O} = \lambda o.\,\mathrm{in}_O\,(\overline{\mathit{node}_{O_l}\,(\pi_l\,(\pi_1\,o))}, \overline{\pi_y\,(\pi_1\,o)}) & \text{its legs, of node types } O_l \text{ smaller than } O\text{, read the same way, its other fields as they are} \\
\llbracket e_1 \texttt{ ? } e_2 \texttt{ : } e_3 \rrbracket = \mathbf{if}\ \llbracket e_1 \rrbracket\ \mathbf{then}\ \llbracket e_2 \rrbracket\ \mathbf{else}\ \llbracket e_3 \rrbracket & \text{case on } \mathit{Bool} \\
\llbracket \iota_1\,e \rrbracket = \mathrm{inl}\,\llbracket e \rrbracket \qquad \llbracket \iota_2\,e \rrbracket = \mathrm{inr}\,\llbracket e \rrbracket & \text{\cppinline{nullopt} is $\iota_1$ of $()$} \\
\llbracket \texttt{case(}s\texttt{, }f\texttt{, }g\texttt{)} \rrbracket = \mathbf{case}\ \llbracket s \rrbracket\ \mathbf{of}\ \mathrm{inl}\,x \Rightarrow \llbracket f \rrbracket^{\!\circ}\,x \mid \mathrm{inr}\,y \Rightarrow \llbracket g \rrbracket^{\!\circ}\,y & \text{case on a sum, each branch elaborated at its summand} \\
\llbracket \texttt{std::make\_shared<const }\tau\texttt{>(}e\texttt{)} \rrbracket = ((), \llbracket e \rrbracket) \qquad \llbracket \texttt{*}e \rrbracket = \pi_2\,\llbracket e \rrbracket & \text{nothing observes which box} \\
\end{array}$}
\]
\normalsize
\caption{The translation.
Types, constraints as record types, dictionaries, declarations and terms; the paragraph \emph{Reading the translation} in Appendix~\ref{app:fjlt:translation} reads it.}
\label{fig:fjlt-translation}
\end{figure}

